\documentclass[reqno,11pt]{amsart}
\usepackage{math327-lecture}
\usepackage{needspace}
\usepackage{tikz}
\lectureinfo{17}{Oct. 19, 2026}{Group Actions and Symmetries of the Square}
\artinsections{6.1, 6.4, 6.7, 6.11}



\tcolorboxenvironment{example}{lecturebox,colback=definitiongreen,colframe=definitionborder}
\definecolor{remarkpurple}{HTML}{F2ECFA}
\definecolor{remarkpurpleborder}{HTML}{A68CC7}
\tcolorboxenvironment{remark}{lecturebox,colback=remarkpurple,colframe=remarkpurpleborder}

\begin{document}
\raggedbottom
\makelecturetitle

\section{A word on group actions}

Let $S$ be a set and $G$ a group. An action of $G$ on $S$ , or a group action, is defined as follows.

\begin{define}[Group action]
A \emph{(left) action} of $G$ on $S$ is a map
\[
  \circ:G\times S\longrightarrow S,\qquad (g,s)\longmapsto g\circ s,
\]
called the \emph{action map}, such that:
\begin{enumerate}[label=\textbf{(\arabic*)}]
  \item $e_G\circ s=s$ for every $s\in S$.
  \item $(gh)\circ s=g\circ(h\circ s)$ for all $g,h\in G$ and $s\in S$.
\end{enumerate}

\end{define}

Let's give some examples which we have already seen before!

\begin{example}[The natural action of $S_n$]
The symmetric group $S_n$ acts on $S=\{1,2,\ldots,n\}$ by evaluation:
\[
  S_n\times S\longrightarrow S,\qquad (\sigma,i)\longmapsto
  \sigma\circ i:=\sigma(i).
\]
The identity permutation fixes every $i\in S$, and composition of
permutations gives
\[
  (\sigma\tau)\circ i=(\sigma\tau)(i)
  =\sigma(\tau(i))=\sigma\circ(\tau\circ i).
\]
Thus both action axioms hold. For example, $(1\,2\,3)\in S_3$ sends
$1$ to $2$, $2$ to $3$, and $3$ to $1$.
\end{example}

\begin{example}[The natural action of $\GL_n(\R)$]
The group $\GL_n(\R)$ of invertible $n\times n$ real matrices acts on
$\R^n$ by matrix multiplication. Regarding vectors as column vectors,
the action map is
\[
  \GL_n(\R)\times\R^n\longrightarrow\R^n,\qquad
  (A,v)\longmapsto A\circ v:=Av.
\]
The identity matrix satisfies $I_n\circ v=I_nv=v$, and associativity
of matrix multiplication gives
\[
  (AB)\circ v=(AB)v=A(Bv)=A\circ(B\circ v).
\]
Thus both action axioms hold.
\end{example}

\begin{example}[The action of $S_4$ on three partitions]
Recall the three partitions from Lecture~16:
\[
\begin{aligned}
  \Pi_1&=\bigl\{\{1,2\},\{3,4\}\bigr\},\\
  \Pi_2&=\bigl\{\{1,3\},\{2,4\}\bigr\},\\
  \Pi_3&=\bigl\{\{1,4\},\{2,3\}\bigr\}.
\end{aligned}
\]
Let $\mathcal P=\{\Pi_1,\Pi_2,\Pi_3\}$. The group $S_4$ acts on
$\mathcal P$ by applying a permutation to every entry:
\[
  \sigma\circ\bigl\{\{a,b\},\{c,d\}\bigr\}
  =\bigl\{\{\sigma(a),\sigma(b)\},
            \{\sigma(c),\sigma(d)\}\bigr\}.
\]
Since $\sigma$ is a bijection, the result is again a partition into
two unordered pairs, so it belongs to $\mathcal P$.
The identity fixes each partition, and applying $\tau$ and then
$\sigma$ to every entry is the same as applying $\sigma\tau$.
Therefore
\[
  (\sigma\tau)\circ\Pi_i=\sigma\circ(\tau\circ\Pi_i),
\]
which verifies the second action axiom.

For example, if $\sigma=(1\,2\,3\,4)$, then
\[
  \sigma\circ\Pi_1=\Pi_3,\qquad
  \sigma\circ\Pi_2=\Pi_2,\qquad
  \sigma\circ\Pi_3=\Pi_1.
\]
Identifying $\Pi_i$ with $i$ gives exactly the homomorphism
$\varphi:S_4\to S_3$ from Lecture~16: if
$\sigma\circ\Pi_i=\Pi_j$, then $\varphi(\sigma)(i)=j$.
In particular, $\varphi((1\,2\,3\,4))=(1\,3)$.
\end{example}


\subsection{Actions as homomorphisms}

The previous example is actually a general fact.

\begin{proposition}
An action of $G$ on a set $S=\{s_1,\ldots,s_d\}$ of cardinality $d$
determines a homomorphism
\[
  \rho:G\longrightarrow S_d.
\]
The correspondence is given by $\rho(g)(i)=j$ if $g\circ s_i=s_j$.
\end{proposition}
\begin{proof}
We need to check that $\rho(g)$ is a permutation of $\{1,\ldots,d\}$
and that the map $g\mapsto\rho(g)$ is a homomorphism, that is,
\[
  \rho(g_1g_2)=\rho(g_1)\rho(g_2).
\]
For the first part, the map $S\to S$ given by $s\mapsto g\circ s$
has inverse $s\mapsto g^{-1}\circ s$. Indeed, if $g\circ s_i=s_j$, then
\[
  g^{-1}\circ s_j=g^{-1}\circ(g\circ s_i)
  =(g^{-1}g)\circ s_i=s_i.
\]
Likewise, $g\circ(g^{-1}\circ s_i)=s_i$.
Thus $\rho(g)$ is a permutation of the indices, with inverse
$\rho(g^{-1})$.

For the homomorphism property, suppose $\rho(g_2)(i)=j$ and
$\rho(g_1)(j)=k$. Then $g_2\circ s_i=s_j$ and $g_1\circ s_j=s_k$, so
\[
  (g_1g_2)\circ s_i=g_1\circ(g_2\circ s_i)=g_1\circ s_j=s_k.
\]
Consequently,
\[
  \rho(g_1g_2)(i)=k=\rho(g_1)(\rho(g_2)(i)).
\]
This holds for every $i$, proving that $\rho$ is a homomorphism.
\end{proof}


\begin{exercise} What is $\ker(\rho)$? Show that \be{} \ker(\rho) = \{ g \in G \mid g \circ s = s \text{ for all } s \in S \}. \ee \end{exercise}

\section{Dihedral group $D_8$}


Label the vertices of a square consecutively as follows:
\begin{center}
\begin{tikzpicture}[scale=1.15]
  \draw[thick] (0,0)--(0,2)--(2,2)--(2,0)--cycle;
  \draw[dashed,courseblue,thick] (-0.35,-0.35)--(2.35,2.35);
  \fill (2,2) circle (1.7pt) node[above right] {$1$};
  \fill (0,2) circle (1.7pt) node[above left] {$2$};
  \fill (0,0) circle (1.7pt) node[below left] {$3$};
  \fill (2,0) circle (1.7pt) node[below right] {$4$};
  \node[courseblue,right] at (2.4,1) {reflection axis};
\end{tikzpicture}
\end{center}
Counterclockwise rotation through $90^\circ$ is the permutation
$x=(1\,2\,3\,4)$, and reflection across the diagonal through vertices
$1$ and $3$ is $r=(2\,4)$. They satisfy
\be{} \label{dihedral-rel}
  x^4=1,\qquad r^2=1,\qquad rxr^{-1}=x^{-1}.
\ee
The last relation expresses the fact that reflecting a rotation reverses
its direction. The picture below follows the vertex labels as we apply
$r^{-1}$, then $x$, then $r$ (products act from right to left).
Since $r^{-1}=r$, both reflections use the dashed diagonal.

\Needspace{13\baselineskip}
\begin{center}
\begin{tikzpicture}[x=1cm,y=1cm,
  every node/.style={font=\small},
  transformarrow/.style={-stealth,thick,courseblue}]
  % Each label moves with its vertex; positions are fixed on the page.
  \foreach \offset/\upperright/\upperleft/\lowerleft/\lowerright in
    {0/1/2/3/4,3.7/1/4/3/2,7.4/2/1/4/3,11.1/2/3/4/1} {
    \begin{scope}[xshift=\offset cm]
      \draw[dashed,ForestGreen] (-0.95,-0.95)--(0.95,0.95);
      \draw[thick] (-0.65,-0.65)--(-0.65,0.65)--(0.65,0.65)
        --(0.65,-0.65)--cycle;
      \foreach \px/\py/\vertex/\anchor in
        {0.65/0.65/\upperright/south west,
         -0.65/0.65/\upperleft/south east,
         -0.65/-0.65/\lowerleft/north east,
         0.65/-0.65/\lowerright/north west} {
        \ifnum\vertex=1
          \fill[courseblue] (\px,\py) circle (2.5pt)
            node[anchor=\anchor,text=courseblue,font=\small\bfseries] {$\vertex$};
        \else
          \fill (\px,\py) circle (1.7pt)
            node[anchor=\anchor] {$\vertex$};
        \fi
      }
    \end{scope}
  }
  \node at (0,1.25) {start};
  \node at (3.7,1.25) {after $r^{-1}$};
  \node at (7.4,1.25) {after $xr^{-1}$};
  \node at (11.1,1.25) {after $rxr^{-1}$};
  \draw[transformarrow] (1.15,0)--(2.55,0)
    node[midway,above] {$r^{-1}=r$}
    node[midway,below,font=\scriptsize] {reflect};
  \draw[transformarrow] (4.85,0)--(6.25,0)
    node[midway,above] {$x$}
    node[midway,below,align=center,font=\scriptsize] {$90^\circ$ counter-\\clockwise};
  \draw[transformarrow] (8.55,0)--(9.95,0)
    node[midway,above] {$r$}
    node[midway,below,font=\scriptsize] {reflect};
  % The same final labeling comes from a clockwise quarter-turn.
  \draw[transformarrow] (0,-1.35)--(0,-1.75)--(11.1,-1.75)--(11.1,-1.35);
  \node[fill=white,inner sep=3pt,text=courseblue] at (5.55,-1.75)
    {$x^{-1}$: one $90^\circ$ clockwise rotation};
\end{tikzpicture}
\end{center}
The final positions of all four labels agree with a clockwise
quarter-turn of the original square, so $rxr^{-1}=x^{-1}$.

If the group $D$ is generated by $x, r$ and satisfies the relations \eqref{dihedral-rel}, then in fact $D$ consist of the following eight elements:

\[
\langle x,r\rangle=\{1,x,x^2,x^3,r,xr,x^2r,x^3r\}.
\]

Now we know that $1, x, x^2, x^3$ are rotations of the square, and $r$ is a reflection. Do $xr, x^2r, x^3r$ also have geometric meanings ? Yes! Let's compute these products from right to left and identify their
reflection axes.

\noindent\begin{minipage}{\textwidth}
\textbf{First example: understanding $xr$.}
Remember that $xr$ means \emph{first reflect by $r=(2\,4)$, then
rotate counterclockwise by $90^\circ$ using $x=(1\,2\,3\,4)$}.
The labels in the pictures move with the vertices. The blue vertex
tracks the original vertex $1$.
\begin{center}
\begin{tikzpicture}[x=1cm,y=1cm,
  every node/.style={font=\small,align=center}]
  \foreach \offset/\panel/\upperright/\upperleft/\lowerleft/\lowerright in
    {0/1/1/2/3/4,4.7/2/1/4/3/2,9.4/3/2/1/4/3} {
    \begin{scope}[xshift=\offset cm]
      \draw[thick] (-0.8,-0.8)--(-0.8,0.8)--(0.8,0.8)
        --(0.8,-0.8)--cycle;
      \ifnum\panel<3
        \draw[dashed,ForestGreen,thick] (-1.05,-1.05)--(1.05,1.05);
      \else
        \draw[dashed,ForestGreen,thick] (0,-1.15)--(0,1.15);
      \fi
      \foreach \px/\py/\vertex/\anchor in
        {0.8/0.8/\upperright/south west,
         -0.8/0.8/\upperleft/south east,
         -0.8/-0.8/\lowerleft/north east,
         0.8/-0.8/\lowerright/north west} {
        \ifnum\vertex=1
          \fill[courseblue] (\px,\py) circle (2.5pt)
            node[anchor=\anchor,text=courseblue,font=\small\bfseries] {$\vertex$};
        \else
          \fill (\px,\py) circle (1.6pt)
            node[anchor=\anchor] {$\vertex$};
        \fi
      }
      \ifnum\panel=1
        \node at (0,1.55) {original square};
        \node[ForestGreen] at (0,-1.5) {axis of $r$};
      \fi
      \ifnum\panel=2
        \node at (0,1.55) {after $r$};
      \fi
      \ifnum\panel=3
        \node at (0,1.55) {after $xr$};
        \node[ForestGreen] at (0,-1.5) {axis of the resulting reflection};
      \fi
    \end{scope}
  }
  \draw[-stealth,courseblue,thick] (1.35,0)--(3.35,0)
    node[midway,above] {$r$}
    node[midway,below,font=\scriptsize] {reflect};
  \draw[-stealth,courseblue,thick] (6.05,0)--(8.05,0)
    node[midway,above] {$x$}
    node[midway,below,align=center,font=\scriptsize]
      {rotate $90^\circ$\\counterclockwise};
\end{tikzpicture}
\end{center}
To compute the product, follow each original vertex through both steps:
\[
\begin{aligned}
  1&\xrightarrow{\ r\ }1\xrightarrow{\ x\ }2,
  &2&\xrightarrow{\ r\ }4\xrightarrow{\ x\ }1,\\
  3&\xrightarrow{\ r\ }3\xrightarrow{\ x\ }4,
  &4&\xrightarrow{\ r\ }2\xrightarrow{\ x\ }3.
\end{aligned}
\]
Thus $xr=(1\,2)(3\,4)$. Comparing the first and last squares, the
left and right vertices have exchanged places while their heights
are unchanged. Therefore $xr$ is reflection across the \emph{vertical
line through the midpoints of the top and bottom edges}.
\end{minipage}

\medskip
\noindent\begin{minipage}{\textwidth}
\textbf{The other two products: $x^2r$ and $x^3r$.}
In each case we still apply $r$ first. We then rotate by $180^\circ$
for $x^2r$, or by $270^\circ$ counterclockwise for $x^3r$.
Following the vertices gives
\[
\begin{array}{c|cccc|c}
  \text{product} & 1\mapsto & 2\mapsto & 3\mapsto & 4\mapsto
    & \text{cycle notation}\\ \hline
  x^2r & 3 & 2 & 1 & 4 & (1\,3)\\
  x^3r & 4 & 3 & 2 & 1 & (1\,4)(2\,3)
\end{array}
\]
The pictures show the resulting reflections on the original labeled
square. Dashed lines are the reflection axes; double arrows indicate
exchanged vertices, and circles mark fixed vertices.
\begin{center}
\begin{tikzpicture}[x=1cm,y=1cm,
  every node/.style={font=\small,align=center}]
  \foreach \offset/\panel/\axiscolor in
    {0/1/courseblue,7/2/BrickRed} {
    \begin{scope}[xshift=\offset cm]
      \draw[thick] (-0.8,-0.8)--(-0.8,0.8)--(0.8,0.8)
        --(0.8,-0.8)--cycle;
      \ifnum\panel=1
        \draw[dashed,\axiscolor,line width=1pt] (-1.1,1.1)--(1.1,-1.1);
        \draw[stealth-stealth,\axiscolor,thick] (-0.6,-0.6)--(0.6,0.6);
        \draw[\axiscolor,thick] (-0.8,0.8) circle (0.1);
        \draw[\axiscolor,thick] (0.8,-0.8) circle (0.1);
        \node[\axiscolor] at (0,1.55) {$x^2r=(1\,3)$};
        \node[\axiscolor] at (0,-1.6) {diagonal through $2$ and $4$};
      \else
        \draw[dashed,\axiscolor,line width=1pt] (-1.25,0)--(1.25,0);
        \draw[stealth-stealth,\axiscolor,thick] (-0.65,-0.6)--(-0.65,0.6);
        \draw[stealth-stealth,\axiscolor,thick] (0.65,-0.6)--(0.65,0.6);
        \node[\axiscolor] at (0,1.55) {$x^3r=(1\,4)(2\,3)$};
        \node[\axiscolor] at (0,-1.6) {horizontal line through edge midpoints};
      \fi
      \fill (0.8,0.8) circle (1.6pt) node[above right] {$1$};
      \fill (-0.8,0.8) circle (1.6pt) node[above left] {$2$};
      \fill (-0.8,-0.8) circle (1.6pt) node[below left] {$3$};
      \fill (0.8,-0.8) circle (1.6pt) node[below right] {$4$};
    \end{scope}
  }
\end{tikzpicture}
\end{center}
For $x^2r$, vertices $2$ and $4$ are fixed, so the reflection axis is
exactly the diagonal joining them; vertices $1$ and $3$ are exchanged.
For $x^3r$, the top and bottom vertices are exchanged while left and
right positions are preserved, giving the horizontal reflection axis.
\end{minipage}

\medskip

So we see that in $D$ we have two types of elements: rotations and reflections.

Note that $D$ acts on $\{ 1, 2, 3, 4\}$ the vertices of the square and this gives us a homomorphism $\rho: D \rightarrow S_4$. What is its kernel? It must be trivial as only the identity element of $D$ fixes all the vertices. So we can really think of $D$ as a subset of $S_4$ if we like.
If we do this,
\[
  x^2=(1\,3)(2\,4),\qquad xr=(1\,2)(3\,4),\qquad
  x^3r=(1\,4)(2\,3),
\]
so this group contains the Klein four group $K$ that came up as the kernel of $\varphi: S_4 \rightarrow S_3.$

You can also see this in pictures:

The three partitions are pictured in the top row. The other rows show
what happens when we apply $x^2$, $xr$, or $x^3r$ to \emph{every entry}
of each partition. Vertex positions and labels are fixed in these
pictures; the colored segments show the resulting pairs.
Blue and red distinguish the two pairs so that we can track them.
\begin{center}
\begin{tikzpicture}[x=1cm,y=1cm,
  every node/.style={font=\small,align=center},
  partitionpair/.style={line width=1.7pt}]
  \node at (-2.05,-0.0) {original};
  \begin{scope}[shift={(0.0,-0.0)}]
    \coordinate (v1) at (0.65,0.65);
    \coordinate (v2) at (-0.65,0.65);
    \coordinate (v3) at (-0.65,-0.65);
    \coordinate (v4) at (0.65,-0.65);
    \draw[gray!55,thin] (v1)--(v2)--(v3)--(v4)--cycle;
    \draw[partitionpair,courseblue] (v1)--(v2);
    \draw[partitionpair,BrickRed] (v3)--(v4);
    \fill (v1) circle (1.6pt) node[above right] {$1$};
    \fill (v2) circle (1.6pt) node[above left] {$2$};
    \fill (v3) circle (1.6pt) node[below left] {$3$};
    \fill (v4) circle (1.6pt) node[below right] {$4$};
    \node at (0,-1.35) {$\bigl\{\textcolor{courseblue}{\{1,2\}},\textcolor{BrickRed}{\{3,4\}}\bigr\}$};
    \node[font=\bfseries] at (0,1.3) {$\Pi_1$};
  \end{scope}
  \begin{scope}[shift={(4.8,-0.0)}]
    \coordinate (v1) at (0.65,0.65);
    \coordinate (v2) at (-0.65,0.65);
    \coordinate (v3) at (-0.65,-0.65);
    \coordinate (v4) at (0.65,-0.65);
    \draw[gray!55,thin] (v1)--(v2)--(v3)--(v4)--cycle;
    \draw[partitionpair,courseblue] (v1)--(v3);
    \draw[partitionpair,BrickRed] (v2)--(v4);
    \fill (v1) circle (1.6pt) node[above right] {$1$};
    \fill (v2) circle (1.6pt) node[above left] {$2$};
    \fill (v3) circle (1.6pt) node[below left] {$3$};
    \fill (v4) circle (1.6pt) node[below right] {$4$};
    \node at (0,-1.35) {$\bigl\{\textcolor{courseblue}{\{1,3\}},\textcolor{BrickRed}{\{2,4\}}\bigr\}$};
    \node[font=\bfseries] at (0,1.3) {$\Pi_2$};
  \end{scope}
  \begin{scope}[shift={(9.6,-0.0)}]
    \coordinate (v1) at (0.65,0.65);
    \coordinate (v2) at (-0.65,0.65);
    \coordinate (v3) at (-0.65,-0.65);
    \coordinate (v4) at (0.65,-0.65);
    \draw[gray!55,thin] (v1)--(v2)--(v3)--(v4)--cycle;
    \draw[partitionpair,courseblue] (v1)--(v4);
    \draw[partitionpair,BrickRed] (v2)--(v3);
    \fill (v1) circle (1.6pt) node[above right] {$1$};
    \fill (v2) circle (1.6pt) node[above left] {$2$};
    \fill (v3) circle (1.6pt) node[below left] {$3$};
    \fill (v4) circle (1.6pt) node[below right] {$4$};
    \node at (0,-1.35) {$\bigl\{\textcolor{courseblue}{\{1,4\}},\textcolor{BrickRed}{\{2,3\}}\bigr\}$};
    \node[font=\bfseries] at (0,1.3) {$\Pi_3$};
  \end{scope}
\end{tikzpicture}
\end{center}

\begin{center}
\begin{tikzpicture}[x=1cm,y=1cm,
  every node/.style={font=\small,align=center},
  partitionpair/.style={line width=1.7pt}]
  \node at (-2.05,-3.25) {$x^2$};
  \begin{scope}[shift={(0.0,-3.25)}]
    \coordinate (v1) at (0.65,0.65);
    \coordinate (v2) at (-0.65,0.65);
    \coordinate (v3) at (-0.65,-0.65);
    \coordinate (v4) at (0.65,-0.65);
    \draw[gray!55,thin] (v1)--(v2)--(v3)--(v4)--cycle;
    \draw[partitionpair,courseblue] (v3)--(v4);
    \draw[partitionpair,BrickRed] (v1)--(v2);
    \fill (v1) circle (1.6pt) node[above right] {$1$};
    \fill (v2) circle (1.6pt) node[above left] {$2$};
    \fill (v3) circle (1.6pt) node[below left] {$3$};
    \fill (v4) circle (1.6pt) node[below right] {$4$};
    \node at (0,-1.35) {$\bigl\{\textcolor{courseblue}{\{3,4\}},\textcolor{BrickRed}{\{1,2\}}\bigr\}$};
    \node at (0,-1.95) {$=\Pi_1$};
  \end{scope}
  \begin{scope}[shift={(4.8,-3.25)}]
    \coordinate (v1) at (0.65,0.65);
    \coordinate (v2) at (-0.65,0.65);
    \coordinate (v3) at (-0.65,-0.65);
    \coordinate (v4) at (0.65,-0.65);
    \draw[gray!55,thin] (v1)--(v2)--(v3)--(v4)--cycle;
    \draw[partitionpair,courseblue] (v3)--(v1);
    \draw[partitionpair,BrickRed] (v4)--(v2);
    \fill (v1) circle (1.6pt) node[above right] {$1$};
    \fill (v2) circle (1.6pt) node[above left] {$2$};
    \fill (v3) circle (1.6pt) node[below left] {$3$};
    \fill (v4) circle (1.6pt) node[below right] {$4$};
    \node at (0,-1.35) {$\bigl\{\textcolor{courseblue}{\{3,1\}},\textcolor{BrickRed}{\{4,2\}}\bigr\}$};
    \node at (0,-1.95) {$=\Pi_2$};
  \end{scope}
  \begin{scope}[shift={(9.6,-3.25)}]
    \coordinate (v1) at (0.65,0.65);
    \coordinate (v2) at (-0.65,0.65);
    \coordinate (v3) at (-0.65,-0.65);
    \coordinate (v4) at (0.65,-0.65);
    \draw[gray!55,thin] (v1)--(v2)--(v3)--(v4)--cycle;
    \draw[partitionpair,courseblue] (v3)--(v2);
    \draw[partitionpair,BrickRed] (v4)--(v1);
    \fill (v1) circle (1.6pt) node[above right] {$1$};
    \fill (v2) circle (1.6pt) node[above left] {$2$};
    \fill (v3) circle (1.6pt) node[below left] {$3$};
    \fill (v4) circle (1.6pt) node[below right] {$4$};
    \node at (0,-1.35) {$\bigl\{\textcolor{courseblue}{\{3,2\}},\textcolor{BrickRed}{\{4,1\}}\bigr\}$};
    \node at (0,-1.95) {$=\Pi_3$};
  \end{scope}
\end{tikzpicture}
\end{center}

\begin{center}
\begin{tikzpicture}[x=1cm,y=1cm,
  every node/.style={font=\small,align=center},
  partitionpair/.style={line width=1.7pt}]
  \node at (-2.05,-6.5) {$xr$};
  \begin{scope}[shift={(0.0,-6.5)}]
    \coordinate (v1) at (0.65,0.65);
    \coordinate (v2) at (-0.65,0.65);
    \coordinate (v3) at (-0.65,-0.65);
    \coordinate (v4) at (0.65,-0.65);
    \draw[gray!55,thin] (v1)--(v2)--(v3)--(v4)--cycle;
    \draw[partitionpair,courseblue] (v2)--(v1);
    \draw[partitionpair,BrickRed] (v4)--(v3);
    \fill (v1) circle (1.6pt) node[above right] {$1$};
    \fill (v2) circle (1.6pt) node[above left] {$2$};
    \fill (v3) circle (1.6pt) node[below left] {$3$};
    \fill (v4) circle (1.6pt) node[below right] {$4$};
    \node at (0,-1.35) {$\bigl\{\textcolor{courseblue}{\{2,1\}},\textcolor{BrickRed}{\{4,3\}}\bigr\}$};
    \node at (0,-1.95) {$=\Pi_1$};
  \end{scope}
  \begin{scope}[shift={(4.8,-6.5)}]
    \coordinate (v1) at (0.65,0.65);
    \coordinate (v2) at (-0.65,0.65);
    \coordinate (v3) at (-0.65,-0.65);
    \coordinate (v4) at (0.65,-0.65);
    \draw[gray!55,thin] (v1)--(v2)--(v3)--(v4)--cycle;
    \draw[partitionpair,courseblue] (v2)--(v4);
    \draw[partitionpair,BrickRed] (v1)--(v3);
    \fill (v1) circle (1.6pt) node[above right] {$1$};
    \fill (v2) circle (1.6pt) node[above left] {$2$};
    \fill (v3) circle (1.6pt) node[below left] {$3$};
    \fill (v4) circle (1.6pt) node[below right] {$4$};
    \node at (0,-1.35) {$\bigl\{\textcolor{courseblue}{\{2,4\}},\textcolor{BrickRed}{\{1,3\}}\bigr\}$};
    \node at (0,-1.95) {$=\Pi_2$};
  \end{scope}
  \begin{scope}[shift={(9.6,-6.5)}]
    \coordinate (v1) at (0.65,0.65);
    \coordinate (v2) at (-0.65,0.65);
    \coordinate (v3) at (-0.65,-0.65);
    \coordinate (v4) at (0.65,-0.65);
    \draw[gray!55,thin] (v1)--(v2)--(v3)--(v4)--cycle;
    \draw[partitionpair,courseblue] (v2)--(v3);
    \draw[partitionpair,BrickRed] (v1)--(v4);
    \fill (v1) circle (1.6pt) node[above right] {$1$};
    \fill (v2) circle (1.6pt) node[above left] {$2$};
    \fill (v3) circle (1.6pt) node[below left] {$3$};
    \fill (v4) circle (1.6pt) node[below right] {$4$};
    \node at (0,-1.35) {$\bigl\{\textcolor{courseblue}{\{2,3\}},\textcolor{BrickRed}{\{1,4\}}\bigr\}$};
    \node at (0,-1.95) {$=\Pi_3$};
  \end{scope}
\end{tikzpicture}
\end{center}

\begin{center}
\begin{tikzpicture}[x=1cm,y=1cm,
  every node/.style={font=\small,align=center},
  partitionpair/.style={line width=1.7pt}]
  \node at (-2.05,-9.75) {$x^3r$};
  \begin{scope}[shift={(0.0,-9.75)}]
    \coordinate (v1) at (0.65,0.65);
    \coordinate (v2) at (-0.65,0.65);
    \coordinate (v3) at (-0.65,-0.65);
    \coordinate (v4) at (0.65,-0.65);
    \draw[gray!55,thin] (v1)--(v2)--(v3)--(v4)--cycle;
    \draw[partitionpair,courseblue] (v4)--(v3);
    \draw[partitionpair,BrickRed] (v2)--(v1);
    \fill (v1) circle (1.6pt) node[above right] {$1$};
    \fill (v2) circle (1.6pt) node[above left] {$2$};
    \fill (v3) circle (1.6pt) node[below left] {$3$};
    \fill (v4) circle (1.6pt) node[below right] {$4$};
    \node at (0,-1.35) {$\bigl\{\textcolor{courseblue}{\{4,3\}},\textcolor{BrickRed}{\{2,1\}}\bigr\}$};
    \node at (0,-1.95) {$=\Pi_1$};
  \end{scope}
  \begin{scope}[shift={(4.8,-9.75)}]
    \coordinate (v1) at (0.65,0.65);
    \coordinate (v2) at (-0.65,0.65);
    \coordinate (v3) at (-0.65,-0.65);
    \coordinate (v4) at (0.65,-0.65);
    \draw[gray!55,thin] (v1)--(v2)--(v3)--(v4)--cycle;
    \draw[partitionpair,courseblue] (v4)--(v2);
    \draw[partitionpair,BrickRed] (v3)--(v1);
    \fill (v1) circle (1.6pt) node[above right] {$1$};
    \fill (v2) circle (1.6pt) node[above left] {$2$};
    \fill (v3) circle (1.6pt) node[below left] {$3$};
    \fill (v4) circle (1.6pt) node[below right] {$4$};
    \node at (0,-1.35) {$\bigl\{\textcolor{courseblue}{\{4,2\}},\textcolor{BrickRed}{\{3,1\}}\bigr\}$};
    \node at (0,-1.95) {$=\Pi_2$};
  \end{scope}
  \begin{scope}[shift={(9.6,-9.75)}]
    \coordinate (v1) at (0.65,0.65);
    \coordinate (v2) at (-0.65,0.65);
    \coordinate (v3) at (-0.65,-0.65);
    \coordinate (v4) at (0.65,-0.65);
    \draw[gray!55,thin] (v1)--(v2)--(v3)--(v4)--cycle;
    \draw[partitionpair,courseblue] (v4)--(v1);
    \draw[partitionpair,BrickRed] (v3)--(v2);
    \fill (v1) circle (1.6pt) node[above right] {$1$};
    \fill (v2) circle (1.6pt) node[above left] {$2$};
    \fill (v3) circle (1.6pt) node[below left] {$3$};
    \fill (v4) circle (1.6pt) node[below right] {$4$};
    \node at (0,-1.35) {$\bigl\{\textcolor{courseblue}{\{4,1\}},\textcolor{BrickRed}{\{3,2\}}\bigr\}$};
    \node at (0,-1.95) {$=\Pi_3$};
  \end{scope}
\end{tikzpicture}
\end{center}
The order within each pair, and the order of the two pairs, do not
matter. For example, $x^2$ sends $\Pi_1$ to
$\{\{3,4\},\{1,2\}\}=\Pi_1$: the two pairs have simply exchanged places.
Every column therefore remains the same partition. Thus
\[
  x^2\circ\Pi_i=xr\circ\Pi_i=x^3r\circ\Pi_i=\Pi_i
  \quad(i=1,2,3),
\]
so all three symmetries lie in $\ker(\varphi)=K$.

\end{document}
